% Vietnamese translation for OI Public Library
% Source-PDF: IOI中国国家候选队论文/2001/刘汝佳(不完整).pdf
\documentclass[11pt]{article}
\usepackage{oipl-vn}

\title{Gợi mở từ bài toán Sokoban}
\author{Liu Rujia}
\date{}

\begin{document}
\maketitle

\origtitle{Insights from the Sokoban problem}
\sourcepath{IOI China candidate team papers / 2001 / Liu Rujia (incomplete).pdf}

\paragraph{Từ khóa.}
Bài toán Sokoban; tìm kiếm trí tuệ nhân tạo; IDA*.

\section*{Tóm tắt}

Bài viết này thảo luận một bài toán vừa thú vị vừa đầy thách thức: bài toán
Sokoban. Bài viết bắt đầu từ các kiến thức cơ bản của tìm kiếm trong không
gian trạng thái, chọn thuật toán IDA* theo đặc điểm của bài toán Sokoban, rồi
thực hiện một số cải tiến bước đầu. Phần chính của bài viết bàn về vài cách
làm cho chương trình thông minh hơn: cải thiện ước lượng cận dưới, nhận biết
bế tắc, phân rã và hợp nhất nhiệm vụ hợp lý, tìm kiếm mẫu và thử nghiệm ngẫu
nhiên. Cuối cùng, bài viết giới thiệu sơ lược một số ý tưởng khác chưa được
nhắc tới ở trước và đưa ra tổng kết.

\renewcommand{\contentsname}{Mục lục}
\tableofcontents

\section*{Lời nói đầu}
\addcontentsline{toc}{section}{Lời nói đầu}

``Sokoban'' là một trò chơi trí tuệ một người rất phổ biến. Nhiệm vụ của
người chơi là điều khiển một công nhân trong nhà kho để đẩy \(N\) thùng giống
nhau tới \(N\) ô đích giống nhau. Nếu chưa rõ luật chơi thì cũng không sao:
chơi thử SokoMind trong phần đính kèm là biết. Tôi đã thêm vào đó 90 màn kiểm
thử chuẩn, 61 màn trẻ em và 170 màn của Wenquxing. Trong phần giới thiệu về
sau, ta sẽ dùng các màn ấy để kiểm thử chương trình, vì vậy tôi khuyên bạn
nên tự thử trước xem mình vượt được bao nhiêu màn.

Vì nội dung bài viết không ít, trước hết tôi đưa ra một gợi ý nhỏ về cách
đọc. Các bạn mới học, hoặc chưa biết thuật toán IDA*, nhất định nên đọc từ
chương 1 và hiểu rõ nó, vì đó là nền tảng của toàn bài. Chương 2 rất dễ hiểu;
chỉ cần nắm đại ý, biết khó khăn của bài toán Sokoban nằm ở đâu và vì sao ta
chọn thuật toán IDA*. Bạn nào nóng lòng muốn xem chương trình có thể đọc
trước chương 3: phiên bản đầu tiên S4-Baby của chúng ta ra đời ở đó, sau đó
là một loạt biện pháp cải tiến. Những bạn thích trí tuệ nhân tạo cũng nên đọc
kỹ, biết đâu sẽ tìm được cảm hứng. Cuối cùng là phần tổng kết, cũng quan
trọng như chương 1, đừng bỏ qua.

Tôi giả định độc giả của bài viết đã có hiểu biết nhất định về các kiến thức
liên quan, nên thường không đưa ra định nghĩa hoặc chứng minh thật chặt chẽ,
mà chỉ nhắc qua một cách thô lược và cố gắng dùng ngôn ngữ dễ hiểu. Nhưng vì
trình độ của tôi thật sự có hạn, chắc chắn có không ít chỗ sai, rất mong mọi
người phê bình và chỉ bảo.

Rất hoan nghênh mọi người liên hệ với tôi. Email: \texttt{liurujia@163.net},
OICQ: \texttt{2575127}.

\section{Kiến thức cơ bản về tìm kiếm không gian trạng thái}

\subsection{Không gian trạng thái}

Đối với một bài toán thực tế, ta có thể trừu tượng hóa nó ở một mức độ nhất
định. Nói một cách thông thường, trạng thái là mô tả toán học của tiến triển
bài toán tại một thời điểm nào đó; chuyển trạng thái là thao tác đưa bài toán
từ một trạng thái sang một, hoặc vài, trạng thái khác. Nếu chỉ có một tác tử
có thể thực hiện các chuyển trạng thái ấy, thì mục tiêu của ta là đơn nhất:
từ trạng thái bắt đầu xác định, qua một chuỗi chuyển trạng thái để đi tới một
hoặc nhiều trạng thái đích.

Nếu có nhiều hơn một tác tử có thể điều khiển chuyển trạng thái, chẳng hạn
trong đánh cờ, thì chúng có thể chuyển trạng thái theo những mục tiêu khác
nhau, thậm chí đối lập nhau. Loại bài toán ấy không thuộc phạm vi thảo luận
của bài viết này.

Ta biết rằng quá trình tìm kiếm thực chất là duyệt một đồ thị ẩn. Các đỉnh
của nó là mọi trạng thái, còn các cạnh có hướng tương ứng với các chuyển
trạng thái. Một lời giải khả thi là một đường đi từ đỉnh bắt đầu tới một đỉnh
bất kỳ trong tập trạng thái đích. Đồ thị ấy gọi là không gian trạng thái; kiểu
tìm kiếm như vậy gọi là tìm kiếm trong không gian trạng thái một tác tử.

\subsection{Tìm kiếm mù}

Tìm kiếm mù chủ yếu gồm các loại sau.

\paragraph{Tìm kiếm thuần ngẫu nhiên.}
Nghe có vẻ khá ``ngốc'', nhưng khi độ sâu rất lớn, số lời giải khả thi nhiều
và độ sâu lời giải lại không quan trọng, nó vẫn hữu dụng. Hơn nữa, tìm kiếm
ngẫu nhiên sau khi cải tiến có thể xử lý các bài toán mà phân bố lời giải có
quy luật nhất định, như tương đối dày đặc, tương đối đều, hoặc phân bố theo
tỉ lệ vàng. Một ví dụ điển hình là khi bạn cuống lên tìm đồ, thường bạn đang
thực hiện tìm kiếm ngẫu nhiên.

\paragraph{Tìm kiếm theo chiều rộng và chiều sâu.}
Mọi người hẳn đã quen thuộc với hiệu suất thời gian, hiệu suất không gian và
đặc điểm của BFS và DFS. Ví dụ của tìm kiếm theo chiều rộng là khi kính của
bạn rơi xuống đất và bạn bò trên sàn để tìm: bạn luôn sờ trước những chỗ gần
mình nhất, nếu không thấy thì sờ xa hơn một chút. Ví dụ điển hình của tìm
kiếm theo chiều sâu là đi trong mê cung. Chúng còn có các hình thức tìm kiếm
ngược và hai chiều, nhưng những phần đó không nằm trong phạm vi thảo luận ở
đây.

\paragraph{Tìm kiếm lặp.}
Các phép tìm kiếm này đặt một số hạn chế lên cách mở rộng cây tìm kiếm, rồi
lặp lại việc tìm kiếm bằng cách nới dần điều kiện. Chúng gồm các phương pháp
sau.

\paragraph{Lặp sâu dần.}
Giới hạn độ sâu tối đa của cây tìm kiếm là \(D_{\max}\), rồi tiến hành tìm
kiếm. Nếu không có lời giải thì tăng \(D_{\max}\) rồi tìm tiếp. Tuy làm nhiều
công việc lặp lại, nhưng vì lượng việc tìm kiếm tăng theo hàm mũ của độ sâu,
lượng việc ở lần trước, tức phần lặp lại, nhỏ hơn nhiều so với lượng tìm kiếm
hiện tại. Phương pháp này thích hợp với các bài toán mà cây tìm kiếm nhìn
chung vừa rộng vừa sâu, nhưng lời giải khả thi lại không quá sâu. DFS thông
thường có thể sa vào một vùng rất sâu mà không có lời giải; còn BFS thì không
đủ không gian.

\paragraph{Lặp rộng dần.}
Phương pháp này giới hạn số con tối đa \(B_{\max}\) được mở rộng từ một đỉnh.
Nhưng vì ưu điểm không thật rõ, nó không được ứng dụng nhiều và cũng ít được
nghiên cứu.

\paragraph{Tìm kiếm chùm.}
Phương pháp này giới hạn tổng số đỉnh tối đa \(W_{\max}\) ở mỗi tầng của cây
tìm kiếm. Rõ ràng khi đó kích thước cây tìm kiếm tỉ lệ thuận với độ sâu,
nhưng có thể bỏ lỡ lời giải rất gần điểm bắt đầu. Khi tăng \(W_{\max}\), việc
giữ lại các đỉnh nào và tăng \(W_{\max}\) bao nhiêu vẫn là vấn đề đang được
nghiên cứu.

\subsection{Tìm kiếm có thông tin}

Ta cảm thấy một số bài toán rất ``có cái để nghĩ'', chủ yếu vì chúng có nhiều
thông tin gợi ý, dễ bắt đầu suy nghĩ, nhưng lại không dễ tìm được lời giải.
Ta không muốn tự tay thử từng trường hợp một cách mù quáng, và cũng không
muốn chương trình tìm kiếm một cách máy móc. Nói cách khác, ta hy vọng khai
thác càng nhiều càng tốt đặc điểm riêng của bài toán để làm cho tìm kiếm
thông minh hơn. Tìm kiếm có thông tin được giới thiệu dưới đây chính là một
dạng tìm kiếm thông minh như vậy.

Trong các thuật toán vừa nhắc tới, ta chưa dùng thông tin của bản thân trạng
thái, mà chỉ dùng chuyển trạng thái để tìm kiếm. Trên thực tế, khi tự giải
bài, ta thường ước lượng xem một trạng thái còn cách mục tiêu bao xa, rồi
dựa vào đó chọn giữa nhiều phương án. Đưa cách làm ấy vào tìm kiếm, ta có
thể dùng một hàm đánh giá của trạng thái để ước lượng khoảng cách từ nó tới
trạng thái đích. Hàm đánh giá này gắn rất chặt với bài toán và thể hiện một
mức độ thông minh nhất định. Để tiện trình bày về sau, ta giới thiệu trước
một số ký hiệu.

\begin{center}
\begin{tabular}{ll}
\(S\) & Một trạng thái bất kỳ của bài toán.\\
\(H^*(s)\) & Khoảng cách thực, ngắn nhất, từ \(s\) tới đích; tiếc là không biết trước.\\
\(H(s)\) & Hàm heuristic của \(s\), tức cận dưới của khoảng cách từ \(s\) tới đích.\\
\(G(s)\) & Chi phí trước khi tới trạng thái \(s\), thường lấy là độ sâu của \(s\).\\
\(F(s)\) & Hàm đánh giá của \(s\), tức ước lượng tổng chi phí để tới đích.\\
\end{tabular}
\end{center}

Với \(H(s)\), ta yêu cầu \(h(s)\le h^*(s)\). Nếu với mọi trạng thái \(s_1\) và
\(s_2\), hàm \(h\) còn thỏa
\[
  h(s_1)\le h(s_2)+c(s_1,s_2),
\]
trong đó \(c(s_1,s_2)\) là chi phí chuyển từ trạng thái \(s_1\) sang \(s_2\),
tức khi chuyển trạng thái thì mức giảm của cận dưới \(h\) nhiều nhất bằng
chi phí thực của chuyển trạng thái, ta nói hàm \(h\) là nhất quán. Thực ra
điều đó chỉ yêu cầu \(h\) không được giảm quá nhanh.

Trực giác cho ta
\[
  f(s)=g(s)+h(s),
\]
tức tổng của chi phí đã trả và chi phí còn phải trả. Nếu \(h\) nhất quán, với
\(s_1\) và một đỉnh hậu duệ \(s_2\) của nó ta có
\[
  h(s_1)\le h(s_2)+c(s_1,s_2).
\]
Cộng \(g(s_1)\) vào hai vế, ta được
\[
  h(s_1)+g(s_1)\le h(s_2)+g(s_1)+c(s_1,s_2),
\]
tức \(f(s_1)\le f(s_2)\). Vì vậy hàm \(f\) tăng đơn điệu.

\paragraph{Tìm kiếm tham lam.}
Giống tìm kiếm theo chiều rộng, nó dùng một hàng đợi để lưu các đỉnh chờ mở
rộng, nhưng sắp theo giá trị hàm \(h\) tăng dần, thực chất là hàng đợi ưu
tiên. Rõ ràng vì \(h\) ước lượng không chính xác, tìm kiếm tham lam không bảo
đảm lời giải đầu tiên tìm được là tối ưu, và có thể rất lâu cũng không tìm
được lời giải nào.

\paragraph{Thuật toán A*.}
Thuật toán này rất giống tìm kiếm tham lam, nhưng sắp theo giá trị hàm \(f\).
Điều đó sinh thêm một vấn đề: trạng thái mới sinh ra có thể đã gặp trước đó.
Vì sao? Tìm kiếm tham lam sắp theo hàm \(h\), mà \(h\) chỉ liên quan tới
trạng thái, nên không xuất hiện trùng lặp kiểu này; còn giá trị \(f\) không
chỉ liên quan tới trạng thái, mà còn liên quan tới cách chuyển tới trạng thái
\(s\), nên cùng một trạng thái có thể có các giá trị \(f\) khác nhau. Cách
giải quyết cũng đơn giản: nếu trạng thái mới \(s_1\) giống một trạng thái
\(s_2\) đã gặp, giữ lại trạng thái có giá trị \(f\) nhỏ hơn. Nếu \(s_2\) là
đỉnh đang chờ mở rộng, vẫn có thể xảy ra \(f(s_2)>f(s_1)\); chỉ các đỉnh đã
mở rộng mới bảo đảm giá trị \(f\) tăng dần. A* bảo đảm tìm được lời giải tối
ưu, nhưng dùng rất nhiều không gian, khó thích ứng với bài toán Sokoban của
chúng ta.

\paragraph{Thuật toán IDA*.}
Vì A* có vấn đề về không gian, liệu ta có thể mượn ưu thế không gian của DFS,
dùng phương thức tìm kiếm lặp để giảm áp lực hay không? Sau nghiên cứu, năm
1985 Korf đưa ra thuật toán Iterative Deepening A*, tức IDA*, giải quyết khá
tốt vấn đề này. Ban đầu, ta đặt giá trị độ sâu tối đa \(D_{\max}\) bằng giá
trị \(h\) của đỉnh bắt đầu, rồi bắt đầu tìm kiếm theo chiều sâu, bỏ qua mọi
đỉnh có giá trị \(f\) lớn hơn \(D_{\max}\), nhờ đó giảm được rất nhiều lượng
tìm kiếm. Nếu không có lời giải, lại tăng \(D_{\max}\) cho tới khi tìm được
một lời giải. Dễ chứng minh lời giải ấy nhất định là tối ưu. Vì đổi sang
hình thức theo chiều sâu, so với A*, IDA* thực dụng hơn:

\begin{enumerate}
  \item Không cần kiểm tra trùng, không cần sắp xếp, chỉ cần dùng ngăn xếp.
  Thao tác đơn giản.
  \item Nhu cầu không gian giảm mạnh, có quan hệ logarit với kích thước cây
  tìm kiếm.
\end{enumerate}

\paragraph{Các tìm kiếm có thông tin khác.}
Những phương pháp này gồm A* kiểu DFS cộng cắt tỉa tối ưu, A* hai chiều, v.v.
Nhưng vì chúng chưa chín muồi hoặc không hữu dụng lắm, ở đây không giới thiệu
thêm. A* còn có một dạng có trọng số; do trong bài toán Sokoban hiệu quả
không rõ, ở đây cũng bỏ qua.

\section{Bài toán Sokoban và đặc điểm của nó}

Sau khi đã hiểu nhất định về các thuật toán tìm kiếm trong không gian trạng
thái, ta xét bài toán Sokoban của mình. Dùng phương pháp nào thì tốt? Trước
hết hãy xem đặc điểm của bài toán.

\subsection{Bài toán Sokoban}

Ở trên ta đã giới thiệu bài toán Sokoban. Ở đây tôi chỉ muốn nêu vài điểm cần
chú ý liên quan tới việc giải bài. Trước hết, kích thước bản đồ Sokoban mà ta
xét tối đa là \(20\times 20\), đủ đáp ứng phần lớn các màn. Để tiện thảo luận
về sau, ta đánh số bản đồ: các cột từ trái sang phải gọi là A, B, C, \ldots,
còn các hàng từ trên xuống gọi là a, b, c, \ldots. Vì việc người đi lại khi
không đẩy thùng là không quan trọng, khi ghi lời giải ta bỏ qua vị trí và di
chuyển của người, chỉ ghi các lần di chuyển của thùng. Hành động của người
có thể suy ra rất dễ từ hành động của thùng.

\paragraph{Ghi chú về hình.}
Tại đây PDF đặt hình màn chuẩn số 1 kèm lời giải. Phần chữ trích xuất bằng
\texttt{pdftotext} không chứa sơ đồ bản đồ đó.

Thế nào, ngay màn đầu đã cần nhiều bước như vậy. Các màn về sau sẽ ngày càng
khó.

\subsection{Đặc điểm của bài toán Sokoban}

Trong lời nói đầu tôi đã nói khá lâu; tôi nghĩ dù trước đây bạn chưa chơi thì
bây giờ cũng đã chơi thử rồi. Cảm giác thế nào? Có phải biến hóa quá nhiều,
khó nắm bắt không? Không chỉ con người khó nắm bắt, ngay cả việc lập trình
cũng trở nên khó hơn rất nhiều. Ta hãy so sánh nó với bài toán 8-puzzle kinh
điển.

\paragraph{1. Bế tắc.}
Người mới chơi rất nhanh sẽ học được bế tắc là gì: chỉ cần đẩy một thùng vào
góc. Rõ ràng bố cục như vậy chơi tiếp cũng hết hy vọng, vì về sau đẩy thế
nào cũng không thể đưa thùng ấy ra khỏi góc. Không ít người chơi đã đúc kết
nhiều kinh nghiệm về bế tắc, nhưng giải quyết vấn đề này một cách hệ thống
không hề dễ. Ta sẽ dùng cả một chương, thực ra cũng không dài lắm, để phân
tích vấn đề này.

\begin{verbatim}
##
#$

Be tac dien hinh: thung bi ket trong goc tuong.
\end{verbatim}

Ta lại nhìn bài toán 8-puzzle. Nó không có bế tắc, vì mỗi bước đều đảo ngược
được. Ở điểm này, Sokoban đau đầu hơn nhiều. Dễ thấy không gian trạng thái
như vậy không phải là đồ thị vô hướng, mà là đồ thị có hướng.

\paragraph{2. Không gian trạng thái.}
Bài toán 8-puzzle mỗi lần nhiều nhất có 4 cách di chuyển, số bước nhiều nhất
cũng chỉ vài chục. Còn Sokoban thì sao? Với một màn khó, một bước có thể có
hơn 100 lựa chọn, toàn bộ lời giải gồm hơn 600 động tác đẩy thùng. Cả hệ số
phân nhánh lẫn độ sâu cây lời giải đều lớn như vậy, không gian trạng thái
đương nhiên không nhỏ.

\paragraph{3. Ước lượng cận dưới.}
Trong tìm kiếm có thông tin, ta cần tính giá trị \(h\), tức cần ước lượng cận
dưới. Bài toán 8-puzzle có nhiều hàm cận dưới tốt, như tổng khoảng cách
``về nhà'', nhưng Sokoban thì sao? Ta không thể trực tiếp tính khoảng cách
``về nhà'', vì chưa rõ nhà của thùng nào là ô nào. Rất tự nhiên, ta có thể
làm một phép ghép tối ưu trên đồ thị hai phía, nhưng cận dưới ấy thế nào?

\paragraph{a. Độ chính xác.}
Với A* và các biến thể của nó, cận dưới càng gần chi phí thực thì nói chung
thuật toán càng hiệu quả. Phép ghép tối ưu của ta chỉ là ``tình huống lý
tưởng''. Trên thực tế, trong rất nhiều trường hợp các thùng ràng buộc lẫn
nhau; việc buộc phải rời tuyến mục tiêu để nhường chỗ cho thùng khác là rất
phổ biến. Chẳng hạn ở màn chuẩn số 50, có thùng phải đi qua ô đích rồi rời
khỏi nó để nhường đường cho thùng khác. Hàm cận dưới của ta trả về 100, nhưng
kết quả tốt nhất hiện có là 370. Chênh lệch thật lớn.

\paragraph{b. Hiệu suất.}
Vì hàm cận dưới được gọi rất thường xuyên, hiệu suất của nó không thể xem
nhẹ. Độ phức tạp tiệm cận thời gian của ghép tối ưu khoảng \(O(N^3)\), lớn
hơn hàm cận dưới của 8-puzzle không biết bao nhiêu lần. Về sau ta sẽ đưa ra
một vài cách cải tiến, nhưng bản chất của nó không thay đổi.

\subsection{Giải bài toán Sokoban như thế nào}

Đã có người chứng minh bài toán Sokoban là NP-hard, xem ra ta vẫn nên xét
tìm kiếm. Nhớ lại tìm kiếm trong không gian trạng thái đã nhắc ở phần trước,
dùng loại nào là tốt?

Vì đây là trò chơi trí tuệ, thông tin gợi ý có thể dùng rất phong phú. Ta
không chỉ cần dùng, mà còn phải dùng càng đầy đủ càng tốt, nên cần dùng tìm
kiếm có thông tin. Phía trên cũng đã nói, không gian trạng thái của Sokoban
rất lớn, A* không còn cách nào, vì vậy ta chọn IDA*: hiện thực đơn giản và
nhu cầu không gian cũng ít.

Nếu Sokoban khó như vậy, vì sao lại có nhiều người giải được khá nhiều màn?
Các màn chuẩn 90 đã bị mọi người hiểu rất kỹ từ nhiều năm trước. Vì con người
thông minh. Họ biết dự đoán, biết sắp xếp, biết học, có trực giác hỗ trợ và
cũng có một tinh thần mạo hiểm nhất định. Họ, trong đó có cả tôi, thường dùng
một số chiến lược giải bài ở ``tầng cao'', vừa hiệu quả vừa linh hoạt.
\texttt{Srbga}: muốn học không? \texttt{Readers}: đương nhiên muốn! Đáng tiếc
các chiến lược ấy không đơn giản để học, cũng không thật có quy luật. Trong
các chương sau, tôi sẽ cố hết sức mô phỏng cách tư duy của con người để thêm
càng nhiều trí thông minh càng tốt cho chương trình của chúng ta.

\section{Dùng thuật toán IDA* giải bài toán Sokoban: hiện thực và cải tiến}

Trong phần trước, ta đã biết IDA* là thuật toán lõi để giải Sokoban. Trong
phần này, ta sẽ dùng IDA* làm một chương trình giải Sokoban. Dù đây là phiên
bản đầu tiên, gọi là S4-Baby, cũng đừng xem thường nó.

\subsection{Khung thuật toán IDA*}

Như đã nói, IDA* là thuật toán A* dựa trên DFS lặp, bỏ qua mọi đỉnh có giá
trị \(f\) lớn hơn giới hạn độ sâu. Vì vậy, không khó viết giả mã khung của
IDA*.

\paragraph{Giả mã 1: Khung thuật toán IDA*.}
\begin{verbatim}
procedure IDA_STAR(StartState)
begin
 PathLimit := H( StartState ) - 1;
 Success := False;
  repeat
     inc(PathLimit);
     StartState.g:= 0;
     Push(OpenStack,StartState);
     repeat
      CurrentState:=Pop(OpenStack);
      If Solution(CurrentState) then
        Success = True
      Else if PathLimit >= CurrentState.g + H(CurrentState) then
        Foreach Child(CurrentState) do
         Push(OpenStack, Child(CurrentState));
     until Success or empty(OpenStack);
  until Success or ResourceLimitsReached;
end;
\end{verbatim}

Đây chỉ là một khung rất thô, chưa làm được việc gì. Nhưng tôi nghĩ mọi người
có thể đang nóng lòng muốn thử sức mạnh của IDA*, vì vậy ta cứ làm một
chương trình cơ bản nhất trước.

\subsection{Chương trình đầu tiên}

Muốn từ khung biến thành chương trình cần điền thêm một số thứ. Ở đây ta lần
lượt thảo luận.

\paragraph{Định dạng tệp vào và ra.}
Tệp vào là một tệp văn bản gồm \(N\) dòng, mỗi dòng là một số ký tự. Ý nghĩa
của các ký tự như sau:

\begin{verbatim}
SPACE  o trong
.      o dich
$      thung
*      thung nam tren o dich
@      nguoi day
+      nguoi day nam tren o dich
#      tuong
\end{verbatim}

Định dạng này nhất quán với định dạng của Xsokoban, SokoMind và Rolling
Stone, nên khá tiện. Dòng đầu của tệp ra là số lần đẩy thùng \(M\); tiếp theo
là \(M\) dòng, mỗi dòng có dạng \texttt{x y direction}, nghĩa là đẩy thùng ở
hàng \(x\), cột \(y\) thêm một bước theo hướng \texttt{direction}. Hướng có
thể là một trong \texttt{left}, \texttt{right}, \texttt{up}, \texttt{down},
với \(1\le x,y\le20\).

\paragraph{Cấu trúc dữ liệu.}
Vì đây là phiên bản đầu tiên, ta không cần nghĩ quá nhiều: chỉ cần chạy được
và lập trình tiện là đủ, tạm thời không xét hiệu suất hay những thứ khác. Tối
ưu là chuyện về sau.

Ta định nghĩa các kiểu dữ liệu mới \texttt{BitString}, \texttt{MazeType},
\texttt{MoveType}, \texttt{StateType} và \texttt{IDAType}. Xem chương trình
trong phụ lục, không khó đoán ý nghĩa và công dụng của chúng. Điều duy nhất
cần giải thích là kiểu \texttt{BitString}. Khi ghi trạng thái, ta xem bản đồ
như một số lớn, mỗi ô là một bit. Khi đó toàn bộ các thùng tạo thành một
\texttt{BitString}; kiểm tra một ô nào đó có thùng, ô đích hay tường hay
không chỉ cần kiểm tra bit tương ứng có bằng 1 hay không. Cách này tuy lãng
phí một ít không gian, nhưng phán đoán nhanh hơn và thao tác cũng đơn giản
hơn.

Ta hợp nhất tọa độ \(x,y\) thành một biến \texttt{position}, trong đó
\[
  \texttt{Position}=(x-1)\cdot \texttt{width}+(y-1).
\]
Ta dùng mảng hằng \texttt{DeltaPos:array[0..3]} để biểu diễn lượng tăng của
\texttt{Position} khi đi lên, xuống, trái, phải.

\paragraph{Thuật toán.}
Để đơn giản, ta thậm chí chưa làm ghép tối ưu, mà dùng tổng khoảng cách từ
mỗi thùng tới ô đích gần nhất làm hàm cận dưới. Nhưng ``khoảng cách'' ở đây
là số lần đẩy. Khi tính bằng hàm \texttt{MinPush}, chỉ cần bỏ qua tất cả các
thùng khác rồi dùng một lần BFS.

\paragraph{Hiệu quả.}
Hiệu quả này thì không nói cũng biết: không vượt được màn chuẩn nào. Nhưng để
chứng minh chương trình là đúng, bạn có thể thử các màn trẻ em, tổng cộng 61
màn.

Gì cơ, màn đầu đã không thấy động tĩnh, sinh ra 180 nghìn đỉnh? Tuy vậy, rất
nhiều màn vẫn qua rất nhanh. Ta lấy 1,000,000 đỉnh làm giới hạn, trên máy
Celeron 300A của tôi chạy hơn mười phút, thu được kết quả kiểm thử sau.

\begin{verbatim}
 No.  buoc dinh     No.  buoc  dinh   No.    buoc    dinh
 1    15 186476  21   8   102   41     11    145
 2    6  24      22   7   110   42     10    118
 3    5  14      23   10  192   43     12    223
 4    6  24      24   10  432   44     8     63
 5    9  31      25   4   23    45     12    138
 6    5  8       26   11  846   46     14    178
 7    6  35      27   3   18    47     8     296
 8    11 39      28   9   38    48     8     156
 9    4  12      29   10  142   49     5     60
 10   5  14      30   8   641   50     11    14451
 11   5  13      31   7   192   51     N/A   >1M
 12   4  19      32   3   12    52     N/A   >1M
 13   4  14      33   11  51    53     8     470
 14   6  20      34   11  332   54     16    24270
 15   6  57      35   16  11118 55     N/A   >1M
 16   12 3947    36   10  242   56     14    3318
 17   6  63      37   9   1171  57     N/A   >1M
 18   11 5108    38   11  556   58     N/A   >1M
 19   10 467     39   10  72    59     11    328
 20   10 1681    40   9   203   60     N/A   >1M
                                61     N/A   >1M
\end{verbatim}

Những màn chưa giải được là 51, 52, 55, 57, 58, 60, 61. Những màn tương đối
khó là 1, 16, 18, 20, 26, 30, 35, 37, 38, 50, 53, 54, 56. Dưới đây ta xem
tình hình ước lượng cận dưới của các ``màn khó'', để thấy cái giá phải trả
cho việc ``lười''.

\begin{verbatim}
man   so buoc toi uu   do sau ban dau   tong dinh   dinh tang tren
1       15           11          186476      7416
16      12           7           3947        844
18      11           10          5108        49
20      10           6           1681        42
26      11           5           846         394
30      8            6           641         200
35      16           3           11118       3464
37      9            4           1171        493
38      11           5           556         250
50      11           6           14451       51
53      8            5           470         48
54      16           9           24270       2562
56      14           4           3318        460
\end{verbatim}

Có thể thấy ước lượng cận dưới liên quan rất lớn tới kích thước cây tìm kiếm.
Hãy nhìn các màn 18, 20, 35, 50, 54, 56. Số đỉnh tầng trên ít biết bao. Nếu
ngay từ đầu đã tìm từ tầng ấy thì... xem ra ta thật sự cần dùng thuật toán
ghép tối ưu.

\subsection{Thử sức mạnh của thuật toán ghép tối ưu}

Được, bây giờ ta dùng thuật toán ghép tối ưu. Ghép tối ưu có thể hiện thực
bằng luồng mạng, nhưng ở đây ta dùng thuật toán nhãn đỉnh đã sửa đổi; tôi
chép chương trình trong sách cho đỡ mất công. Bây giờ chương trình đổi tên
thành Baby2.

Sau đây là tình hình kiểm thử của các ``bài khó'' vừa rồi.

\begin{verbatim}
man   buoc thuc   do sau dau   tong dinh Baby-1   tong dinh Baby-2
1     15          15           186476             60
16    12          10           3947               304
18    11          11           5108               46
20    10          8            1681               76
26    11          5            846                552
30    8           8            641                153
35    16          4            11118              6504
37    9           5            1171               438
38    11          5            556                546
50    11          7            14451              98
53    8           8            470                37
54    16          12           24270              273
56    14          4            3318               2225
\end{verbatim}

Ồ, có trường hợp còn ít hơn cả số đỉnh tầng trên vừa rồi. Tất nhiên, cận
dưới tốt hơn thì cắt tỉa theo độ sâu của tầng hiện tại cũng chính xác hơn.

Ngoài ra, bây giờ ta xem 12 màn đầu của Wenquxing. Các màn 1, 2, 4, 6, 8, 9,
11 đã có thể ra lời giải trong 50,000 đỉnh.

\begin{verbatim}
man      buoc thuc  do sau dau   tong dinh   dinh tang tren
1        31         31           75          75
2        11         11           142         142
4        26         18           33923       159
6        16         16           47          47
8        27         21           239         213
9        12         6            4806        2778
11       14         14           73          73
\end{verbatim}

Vậy bước tiếp theo làm gì? In từng trạng thái ra phân tích, ta không khó phát
hiện có rất nhiều đỉnh lặp. Vì vậy vấn đề tự nhiên tiếp theo là: làm sao
tránh lặp?

\subsection{Thử sức mạnh của bảng băm}

Muốn kiểm tra trùng thì đương nhiên cần bảng băm. Nhưng có một vấn đề rất
khó xử lý: biểu diễn trạng thái thế nào? Trong mở rộng đỉnh, ta dùng dòng bit
để định nghĩa trạng thái của các thùng, nhưng ở đây thứ ta cần là chỉ số mảng
phù hợp. Cách biểu diễn đó không thuận tiện. Vì vậy khi xây dựng bảng băm, ta
dùng tọa độ các thùng để biểu diễn trạng thái, tức bộ \(N\) phần tử
\((p[1],p[2],p[3],\ldots,p[n])\). Còn hàm băm thì ta xét theo số mục của bảng
băm. Ở đây, nếu số thùng nhiều nhất là 100, ta thử dùng 10,000 mục. Một
phương án là cộng tất cả tọa độ lại, nhưng cách này có rất nhiều xung đột:
vì khi thùng \(A\) dịch sang phải một ô, thùng \(B\) dịch sang trái một ô,
giá trị hàm băm không đổi. Xét rằng cần giảm xung đột và hàm không nên quá
phức tạp, ta dùng tổng có trọng số của tọa độ làm giá trị băm, tức
\(\sum k\cdot\texttt{Position}[k]\), dĩ nhiên cuối cùng lấy modulo 10,000.
Thực ra hàm này cũng không tốt, nhưng tôi hơi lười, để sau cải tiến. Về xử
lý xung đột, để đơn giản ta dùng phương pháp nối dây, lập danh sách liên kết
để lưu mọi phần tử. Cần chú ý rằng hai trạng thái có cùng tọa độ thùng nhưng
tọa độ người không liên thông là hai trạng thái khác nhau, đều phải lưu. Sau
đây là kết quả kiểm thử các màn vừa rồi.

\begin{verbatim}
man       buoc thuc   do sau dau   dinh Baby2   dinh Baby3
Kid 1     15          15           60           52
Kid 16    12          10           304          189
Kid 18    11          11           46           41
Kid 20    10          8            76           67
Kid 26    11          5            552          192
Kid 30    8           8            153          145
Kid 35    16          4            6504         704
Kid 37    9           5            438          136
Kid 38    11          5            546          152
Kid 50    11          7            98           96
Kid 53    8           8            37           24
Kid 54    16          12           273          258
Kid 56    14          4            2225         1518
Wqx 1     31          31           75           75
Wqx 2     11          11           142          107
Wqx 4     26          18           33923        33916
Wqx 6     16          16           47           46
Wqx 8     27          21           239          226
Wqx 9     12          6            4806         968
Wqx 11    14          14           73           67
\end{verbatim}

Các màn mới hoàn thành là:

\begin{verbatim}
man       buoc thuc   do sau dau   tong dinh   dinh tang tren
Kid 51    13          13           629         629
Kid 52    18          18           39841       39841
Kid 55    14          4            4886        1919
Kid 60    15          9            6916        916
\end{verbatim}

Cần chú ý rằng khi chạy \texttt{kid57} trong chế độ bảo vệ đã xuất hiện tràn
heap, cho thấy việc lưu toàn bộ đỉnh là không cần thiết, vừa tốn không gian,
vừa tốn thời gian, và cũng không mấy khả thi. Vậy ta nên làm gì? Ta biết đánh
giá một bảng băm tốt hay xấu thường xét từ hai mặt: tần suất tra bảng thành
công và lượng công việc tiết kiệm được sau khi tra thành công. Vì vậy ta có
thể đặt hai bảng băm: một bảng lưu các đỉnh gần đây, khả năng tra trúng lớn;
một bảng lưu các đỉnh có độ sâu lớn, vì một khi tìm thấy thì tiết kiệm rất
nhiều công việc. Làm như vậy không tăng bao nhiêu số đỉnh, nhưng giúp hiệu
suất chương trình cao hơn và năng lực giải bài, tức khả năng chịu đựng về
không gian, cũng cải thiện đáng kể. Tuy nhiên để tiện, chương trình của ta
tạm thời chỉ dùng bảng thứ nhất.

\subsection{Tối ưu thứ tự mở rộng đỉnh}

Cải tiến cuối cùng trong phần này là tối ưu thứ tự mở rộng đỉnh. Mục tiêu
không phải là cắt tỉa cây tìm kiếm, mà là hy vọng tìm được lời giải sớm hơn.
Cách cải tiến cụ thể như sau:

\begin{enumerate}
  \item Ưu tiên đẩy thùng vừa được đẩy.
  \item Sau đó thử mọi phương án có thể làm giảm cận dưới; phương án giảm
  nhiều hơn được thử trước. Nếu mức giảm như nhau, đẩy thùng gần đích nhất
  trước.
  \item Cuối cùng thử các phương án khác, cũng xét theo thứ tự như mục 2.
\end{enumerate}

Có thể dự đoán rằng sau khi xử lý như vậy, việc tìm lời giải ``tương đối dễ''
sẽ xảy ra sớm hơn, nhưng cái giá do ước lượng cận dưới không chính xác gây ra
không thể giảm, tức chỉ có thể giảm số đỉnh tầng trên. Tuy vậy, với tư cách
một phương pháp cải tiến chuẩn của IDA*, ta vẫn cần thêm nó vào chương trình
để thử.

Cần chú ý rằng ta dùng ngăn xếp, nên các phương án tương đối kém cần được đưa
vào ngăn xếp trước. Kết quả kiểm thử thực tế cho thấy mục 1 hiệu quả khá tốt,
còn mục 2 và 3 hiệu quả không tốt, thậm chí sinh nhiều đỉnh hơn. Nguyên nhân
chủ yếu có thể là ước lượng cận dưới của ta không chính xác, trong khi mục 2
và 3 dùng tới hàm cận dưới. Trong phiên bản Baby-4 này, ta tắt hai biện pháp
ở mục 2 và 3.

Được rồi, đã viết bốn phiên bản Baby, bạn có muốn so sánh không? Nhưng tôi
chỉ quan tâm vài dữ liệu hơi khó.

\begin{verbatim}
man       buoc thuc   Baby-1     Baby-2     Baby-3   Baby-4
Kid 1     11          186476     60         52       38
Kid 16    7           3947       304        189      149
Kid 18    10          5108       46         41       31
Kid 35    16          11118      6504       704      462
Kid 50    11          14451      98         96       152
Kid 51    13          Too many   Too many   629      54
Kid 52    18          Too many   Too many   39841    97
Kid 54    16          24270      273        258      140
Kid 55    14          Too many   Too many   4886     3390
Kid 56    14          3318       2225       1518     1069
Kid 60    15          Too many   Too many   6916     5022
Wqx 4     26          97855      33923      33916    24251
Wqx 9     12          116927     4806       968      350
\end{verbatim}

Từ bảng trên có thể thấy tối ưu của ta nhìn chung là có hiệu quả. Trực quan
mà nói, nhiều trường hợp cải tiến không rõ là do ước lượng cận dưới tương đối
kém; điểm này về sau ta sẽ tiếp tục thảo luận. Dù thế nào, vượt được 58 trong
61 màn ``trẻ em'' cũng khá tốt, ít nhất có thể nói phiên bản Baby của chương
trình đã có ``trí tuệ'' của một trẻ em bình thường. Nhưng đây chỉ là mở đầu,
phần hay còn ở phía sau.

\subsection{Mã nguồn Baby-4}

Chương trình \texttt{S4BABY4.PAS} nằm trong phần đính kèm; ở đây chỉ thêm một
lượng nhỏ chú thích. Mọi người có thể thử hiệu quả của nó, nhưng không cần
xem quá kỹ, vì trong các chương sau tôi sẽ sửa rất nhiều thứ, thậm chí khung
chương trình chính của IDA* cũng sẽ trở nên khác.

\paragraph{Định nghĩa hằng.}
\begin{verbatim}
const
 {Version}
 VerStr='S4 - SRbGa Super Sokoban Solver (Baby Version 4)';
 Author='Written by Liu Rujia(SrbGa), 2001.2, Chongqing, China';

 {Files}
 InFile='soko.in';
 OutFile='soko.out';

 {Charactors}
 Char_Soko='@';
 Char_SokoInTarget='+';
 Char_Box='$';
 Char_BoxInTarget='*';
 Char_Target='.';
 Char_Wall='#';
 Char_Empty=' ';

 {Dimentions}
 Maxx=21;
 Maxy=21;
 MaxBox=50;

 {Directions}
 Up=0;
 Down=1;
 Left=2;
 Right=3;
 DirectionWords:array[0..3] of string=('UP','DOWN','LEFT','RIGHT');

 {Movement}
 MaxPosition:integer=Maxx*Maxy;
 Opposite:array[0..3] of integer=(1,0,3,2);
 DeltaPos:array[0..3] of integer=(-Maxy,Maxy,-1,1);
\end{verbatim}

Ta hợp nhất tọa độ \(x,y\) thành một giá trị \texttt{position}, trong đó
\(\texttt{position}=(x-1)\cdot\texttt{maxy}+(y-1)\). Ở đây dùng hằng có kiểu
vì về sau sẽ thay đổi giá trị \texttt{MaxPosition} theo kích thước bản đồ.
\texttt{Opposite} là hướng ngược, ví dụ \texttt{Opposite[UP]:=DOWN};
\texttt{DeltaPos} cũng sẽ được đặt lại. Khi di chuyển, ta chỉ cần dùng
\texttt{NewPos:=OldPos+DeltaPos[Direction]}, rất tiện.

\begin{verbatim}
 {IDA Related}
 MaxNode=1000000;
 MaxDepth=100;
 MaxStack=150;
 DispNode=1000;
\end{verbatim}

Cứ sinh bao nhiêu đỉnh thì báo cáo một lần.

\begin{verbatim}
 {HashTable}
 MaxHashEntry=10000;
 HashMask=10000;
 MaxSubEntry=100;

 {BitString}
 BitMask:array[0..7] of byte=(1,2,4,8,16,32,64,128);

 Infinite=Maxint;
\end{verbatim}

\paragraph{Định nghĩa kiểu.}
\begin{verbatim}
type
 PositionType=integer;

 BitString=array[0..Maxx*Maxy div 8-1] of byte;
\end{verbatim}

Toàn bộ bản đồ là một \texttt{BitString}. Bit thứ \texttt{position} bằng 1
khi và chỉ khi vị trí \texttt{position} có thứ gì đó, như thùng, đích hoặc
tường.

\begin{verbatim}
 MapType=array[1..Maxx] of string[Maxy];
 BiGraph=array[1..MaxBox,1..MaxBox] of integer;

 MazeType=
 record
  X,Y:integer;
  Map:MapType;
  GoalPosition:array[1..MaxBox] of integer;
  BoxCount:integer;
  Goals:BitString;
  Walls:BitString;
 end;
\end{verbatim}

Các trường lần lượt là kích thước, dữ liệu gốc dùng để hiển thị trạng thái,
\texttt{BitString} của các ô đích, tổng số thùng, vị trí đích, và
\texttt{BitString} của tường. Cả \texttt{BitString} lẫn mảng vị trí đều được
dùng để tăng tốc.

\begin{verbatim}
 MoveType=
 record
  Position:integer;
  Direction:0..3;
 end;
\end{verbatim}

\texttt{Direction} là hướng mà thùng bị đẩy tới.

\begin{verbatim}
 StateType=
 record
  Boxes:BitString;
  ManPosition:PositionType;
  MoveCount:integer;
  Move:array[1..MaxDepth] of MoveType;
  g,h:integer;
 end;

 IDAType=
 record
  TopLevelNodeCount:longint;
  NodeCount:longint;
  StartState:StateType;
  PathLimit:integer;
  Top:integer;
  Stack:array[1..MaxStack] of StateType;
 end;
\end{verbatim}

\texttt{Top} là con trỏ đỉnh ngăn xếp.

\begin{verbatim}
 PHashTableEntry=^HashTableEntry;
 HashTableEntry=
 record
  Next:PHashTableEntry;
  State:StateType;
 end;

 PHashTableType=^HashTableType;
 HashTableType=
 record
  FirstEntry:array[0..MaxHashEntry] of PHashTableEntry;
  Count:array[0..MaxHashEntry] of byte;
 end;
\end{verbatim}

Đây là các kiểu liên quan tới bảng băm. Ta dùng phương pháp nối dây, nhờ đó
có thể xin không gian heap bằng con trỏ; kết hợp với chế độ bảo vệ thì hiệu
quả tốt hơn.

\begin{verbatim}
var
 HashTable:PHashTableType;
 SokoMaze:MazeType;
 IDA:IDAType;

procedure SetBit(var BS:BitString; p:integer);
begin
 BS[p div 8]:=BS[p div 8] or BitMask[p mod 8];
end;

procedure ClearBit(var BS:BitString; p:integer);
begin
 BS[p div 8]:=BS[p div 8] xor BitMask[p mod 8];
end;

function GetBit(var BS:BitString; p:integer):byte;
begin
 if BS[p div 8] and BitMask[p mod 8]>0 then GetBit:=1 else GetBit:=0;
end;
\end{verbatim}

Đây là các thao tác bit: đặt, xóa và lấy một mục của \texttt{BitString}.

\begin{verbatim}
procedure Init;
var
 Lines:MapType;

 procedure ReadInputFile;
 var
  f:text;
  s:string;
 begin
  SokoMaze.X:=0;
  SokoMaze.Y:=0;
  SokoMaze.BoxCount:=0;
  assign(f,infile);
  reset(f);
  while not eof(f) do
  begin
    readln(f,s);
    if length(s)>SokoMaze.Y then
      SokoMaze.Y:=length(s);
    inc(SokoMaze.X);
    Lines[SokoMaze.X]:=s;
  end;
  close(f);
 end;

 procedure AdjustData;
 var
  i,j:integer;
 begin
  for i:=1 to SokoMaze.X do
    while length(Lines[i])<SokoMaze.Y do
      Lines[i]:=Lines[i]+' ';

 SokoMaze.Map:=Lines;
 for i:=1 to SokoMaze.X do
  for j:=1 to SokoMaze.Y do
    if SokoMaze.Map[i,j] in [Char_BoxInTarget,Char_SokoInTarget,Char_Target] then
      SokoMaze.Map[i,j]:=Char_Target
   else if SokoMaze.Map[i,j]<>Char_Wall then
      SokoMaze.Map[i,j]:=Char_Empty;
\end{verbatim}

Điều chỉnh mảng \texttt{Map}, bỏ thùng và người đẩy ra.

\begin{verbatim}
  for i:=1 to SokoMaze.X do
   for j:=1 to SokoMaze.Y do
    if Lines[i,j] in [Char_Target,Char_BoxInTarget,Char_SokoInTarget] then
    begin
      inc(SokoMaze.BoxCount);
      SokoMaze.GoalPosition[SokoMaze.BoxCount]:=(i-1)*SokoMaze.Y+j-1;
    end;
\end{verbatim}

Thống kê số ô đích và \texttt{GoalPosition}.

\begin{verbatim}
  DeltaPos[Up]:=-SokoMaze.Y;
  DeltaPos[Down]:=SokoMaze.Y;
  MaxPosition:=SokoMaze.X*SokoMaze.Y;
\end{verbatim}

Điều chỉnh \texttt{DeltaPos} và \texttt{MaxPosition} theo kích thước bản đồ.

\begin{verbatim}
 end;

 procedure ConstructMaze;
 var
  i,j:integer;
 begin
  fillchar(SokoMaze.Goals,sizeof(SokoMaze.Goals),0);
  fillchar(SokoMaze.Walls,sizeof(SokoMaze.Walls),0);

  for i:=1 to SokoMaze.X do
   for j:=1 to SokoMaze.Y do
    case Lines[i,j] of
      Char_SokoInTarget, Char_BoxInTarget, Char_Target:
        SetBit(SokoMaze.Goals,(i-1)*SokoMaze.Y+j-1);
      Char_Wall:
        SetBit(SokoMaze.Walls,(i-1)*SokoMaze.Y+j-1);
    end;
 end;

 procedure InitIDA;
 var
  i,j:integer;
  StartState:StateType;
 begin
  IDA.NodeCount:=0;
  IDA.TopLevelNodeCount:=0;
  fillchar(StartState,sizeof(StartState),0);

  for i:=1 to SokoMaze.X do
   for j:=1 to SokoMaze.Y do
    case Lines[i,j] of
      Char_Soko, Char_SokoInTarget:
        StartState.ManPosition:=(i-1)*SokoMaze.Y+j-1;
      Char_Box, Char_BoxInTarget:
        SetBit(StartState.Boxes,(i-1)*SokoMaze.Y+j-1);
    end;

  StartState.g:=0;
  IDA.StartState:=StartState;
  new(HashTable);
  for i:=1 to MaxHashEntry do
  begin
   HashTable^.FirstEntry[i]:=nil;
   HashTable^.Count[i]:=0;
  end;
 end;

begin
 ReadInputFile;
 AdjustData;
 ConstructMaze;
 InitIDA;
end;

procedure PrintState(State:StateType);
var
 i,x,y:integer;
 Map:MapType;
begin
 Map:=SokoMaze.Map;
 x:=State.ManPosition div SokoMaze.Y+1;
 y:=State.ManPosition mod SokoMaze.Y+1;
 if Map[x,y]=Char_Target then
   Map[x,y]:=Char_SokoInTarget
 else
   Map[x,y]:=Char_Soko;

 for i:=1 to MaxPosition do
  if GetBit(State.Boxes,i)>0 then
  begin
    x:=i div SokoMaze.Y+1;
    y:=i mod SokoMaze.Y+1;
    if Map[x,y]=Char_Target then
      Map[x,y]:=Char_BoxInTarget
    else
      Map[x,y]:=Char_Box;
  end;
 for i:=1 to SokoMaze.X do
  Writeln(Map[i]);
end;

function Solution(State:StateType):boolean;
var
 i:integer;
begin
 Solution:=false;
 for i:=1 to MaxPosition do
   if (GetBit(State.Boxes,i)>0) and (GetBit(SokoMaze.Goals,i)=0) then
     exit;
 Solution:=true;
end;

function CanReach(State:StateType; Position:integer):boolean;
\end{verbatim}

Dùng BFS để xét trong trạng thái \texttt{State}, người đẩy có thể tới
\texttt{Position} hay không.

\begin{verbatim}
var
 Direction:integer;
 Pos,NewPos:integer;
 Get,Put:integer;
 Queue:array[0..Maxx*Maxy] of integer;
 Reached:Array[0..Maxx*Maxy] of boolean;
begin
 fillchar(Reached,sizeof(Reached),0);

 Pos:=State.ManPosition;
 Get:=0; Put:=1;
 Queue[0]:=Pos;
 Reached[Pos]:=true;

 CanReach:=true;
 while Get<>Put do
 begin
  Pos:=Queue[Get];
  inc(Get);
  if Pos=Position then
    exit;
  for Direction:=0 to 3 do
  begin
    NewPos:=Pos+DeltaPos[Direction];
    if Reached[NewPos] then continue;
    if GetBit(State.Boxes,NewPos)>0 then continue;
    if GetBit(SokoMaze.Walls,NewPos)>0 then continue;
    Reached[NewPos]:=true;
    Queue[Put]:=NewPos;
    inc(Put);
  end;
 end;
 CanReach:=false;
end;

function MinPush(BoxPosition,GoalPosition:integer):integer;
\end{verbatim}

Khi không có thùng khác, từ \texttt{BoxPosition} đẩy tới \texttt{GoalPosition}
cần ít nhất bao nhiêu bước.

\begin{verbatim}
var
 i:integer;
 Direction:integer;
 Pos,NewPos,ManPos:integer;
 Get,Put:integer;
 Queue:array[0..Maxx*Maxy] of integer;
 Distance:Array[0..Maxx*Maxy] of integer;
begin
 for i:=0 to Maxx*Maxy do
   Distance[i]:=Infinite;

 Pos:=BoxPosition;
 Get:=0; Put:=1;
 Queue[0]:=Pos;
 Distance[Pos]:=0;

 while Get<>Put do
 begin
  Pos:=Queue[Get];
  inc(Get);
  if Pos=GoalPosition then
  begin
    MinPush:=Distance[Pos];
    exit;
  end;
  for Direction:=0 to 3 do
  begin
    NewPos:=Pos+DeltaPos[Direction];
    ManPos:=Pos+DeltaPos[Opposite[Direction]];
    {Nguoi phai dung o phia sau}
    if Distance[NewPos]<Infinite then continue;
    if GetBit(SokoMaze.Walls,NewPos)>0 then continue;
    {Khong day duoc}
    if GetBit(SokoMaze.Walls,ManPos)>0 then continue;
    {Nguoi khong co cho dung}
    Distance[NewPos]:=Distance[Pos]+1;
    Queue[Put]:=NewPos;
    inc(Put);
  end;
 end;
 MinPush:=Infinite;
end;

procedure DoMove(State:StateType; Position,Direction:integer;
                 var NewState:StateType);
var
 NewPos:integer;
begin
 NewState:=State;
 NewPos:=Position+DeltaPos[Direction];
 NewState.ManPosition:=Position;
 SetBit(NewState.Boxes,NewPos);
 ClearBit(NewState.Boxes,Position);
end;

function MinMatch(BoxCount:integer;Gr:BiGraph):integer;
\end{verbatim}

Đây là thuật toán chuẩn, chép từ chương trình trong sách, không cần xem kỹ.

\begin{verbatim}
var
 VeryBig:integer;
 TempGr:BiGraph;
 L:array[1..MaxBox*2] of integer;
 SetX,SetY,MatchedX,MatchedY:Set of 1..MaxBox;

procedure MaxMatch(n,m:integer);
function Path(x:integer):boolean;
var
 i,j:integer;
begin
 Path:=false;
 for i:=1 to m do
   if not (i in SetY)and(Gr[x,i]<>0) then
   begin
     SetY:=SetY+[i];
     if not (i in MatchedY) then
     begin
       Gr[x,i]:=-Gr[x,i];
       MatchedY:=MatchedY+[i];
       Path:=true;
       exit;
     end;
     j:=1;
     while (j<=m)and not (j in SetX) and (Gr[j,i]>=0) do inc(j);
     if j<=m then
     begin
       SetX:=SetX+[j];
       if Path(j) then
       begin
         Gr[x,i]:=-Gr[x,i];
         Gr[j,i]:=-Gr[j,i];
         Path:=true;
         exit;
       end;
     end;
   end;
end;

var
 u,i,j,al:integer;

begin
 Fillchar(L,sizeof(L),0);
 TempGr:=Gr;
 for i:=1 to n do
  for j:=1 to m do
    if L[i]<Gr[i,j] then
      L[i]:=Gr[i,j];
 u:=1; MatchedX:=[]; MatchedY:=[];
 for i:=1 to n do
  for j:=1 to m do
    if L[i]+L[n+j]=TempGr[i,j] then
      Gr[i,j]:=1
    else
      Gr[i,j]:=0;
 while u<=n do
 begin
  SetX:=[u]; SetY:=[];
  if not (u in MatchedX) then
  begin
    if not Path(u) then
    begin
      al:=Infinite;
      for i:=1 to n do
       for j:=1 to m do
         if (i in SetX) and not (j in SetY) and (L[i]+L[n+j]-TempGr[i,j]<al) then
           al:=L[i]+L[n+j]-TempGr[i,j];
      for i:=1 to n do if i in SetX then L[i]:=L[i]-al;
      for i:=1 to m do if i in SetY then l[n+i]:=l[n+i]+al;
      for i:=1 to n do
       for j:=1 to m do
         if l[i]+l[n+j]=TempGr[i,j] then
           Gr[i,j]:=1
         else
           Gr[i,j]:=0;
      MatchedX:=[]; MatchedY:=[];
      for i:=1 to n+m do
       if l[i]<-1000 then
         exit;
    end
    else
      MatchedX:=MatchedX+[u];
    u:=0;
  end;
  inc(u);
 end;
end;

var
 i,j:integer;
 Tot:integer;
begin
 VeryBig:=0;
 for i:=1 to BoxCount do
  for j:=1 to BoxCount do
    if (Gr[i,j]<Infinite)and(Gr[i,j]>VeryBig) then
      VeryBig:=Gr[i,j];
 inc(VeryBig);

 for i:=1 to BoxCount do
  for j:=1 to BoxCount do
    if Gr[i,j]<Infinite then
      Gr[i,j]:=VeryBig-Gr[i,j]
    else
      Gr[i,j]:=0;
\end{verbatim}

Các câu lệnh này thực hiện phép chuyển sang bài toán bổ sung.

\begin{verbatim}
 MaxMatch(BoxCount,BoxCount);
 Tot:=0;
 for i:=1 to BoxCount do
 begin
  for j:=1 to BoxCount do
    if Gr[i,j]<0 then
    begin
      Tot:=Tot+VeryBig-TempGr[i,j];
      break;
    end;
  if Gr[i,j]>=0 then
  begin
    MinMatch:=Infinite;
    exit;
  end;
 end;
 MinMatch:=Tot;
end;

function CalcHeuristicFunction(State:StateType):integer;
\end{verbatim}

Tính giá trị hàm heuristic.

\begin{verbatim}
var
 H,Min:integer;
 i,j,p,Count,BoxCount,Cost:integer;
 BoxPos:array[1..MaxBox] of integer;
 Distance:BiGraph;
begin
 p:=0;
 for i:=1 to MaxPosition do
   if GetBit(State.Boxes,i)>0 then
  begin
    inc(p);
    BoxPos[p]:=i;
  end;
 for i:=1 to p do
  for j:=1 to p do
    Distance[i,j]:=MinPush(BoxPos[i],SokoMaze.GoalPosition[j]);
 BoxCount:=SokoMaze.BoxCount;

 H:=0;
 for i:=1 to BoxCount do
 begin
  Count:=0;
  for j:=1 to BoxCount do
    if Distance[i,j]<Infinite then
      inc(Count);
  if Count=0 then
  {Co mot thung khong day toi bat ky dich nao}
  begin
    CalcHeuristicFunction:=Infinite;
    exit;
  end;
 end;

 H:=MinMatch(BoxCount, Distance);
 CalcHeuristicFunction:=H;
end;

function HashFunction(State:StateType):integer;
var
 i,h,p:integer;
begin
 h:=0;
 p:=0;
 for i:=1 to MaxPosition do
   if GetBit(State.Boxes,i)>0 then
   begin
     inc(p);
     h:=(h+p*i) mod HashMask;
     {Ban co the tu doi ham nay}
   end;
 HashFunction:=h;
end;

function SameState(S1,S2:StateType):boolean;
var
 i:integer;
begin
 SameState:=false;
 for i:=1 to MaxPosition do
   if GetBit(S1.Boxes,i)<>GetBit(S2.Boxes,i) then
     exit;
 if not CanReach(S1,S2.ManPosition) then
 {Chi can hai vi tri nguoi lien thong thi nen xem la cung mot trang thai}
   exit;
 SameState:=true;
end;

function Prior(State:StateType;M1,M2:MoveType):boolean;
var
 NewPos:integer;
 Inertia1,Inertia2:boolean;
 S1,S2:StateType;
 H1,H2:integer;
begin
 Prior:=false;
 if State.MoveCount>0 then
 begin
   NewPos:=State.Move[State.MoveCount].Position+
         DeltaPos[State.Move[State.MoveCount].Direction];
   if NewPos=M1.Position then Inertia1:=true else Inertia1:=false;
   {Uu tien thao tac day lien tiep cung mot thung}
   if NewPos=M2.Position then Inertia2:=true else Inertia2:=false;
   if Inertia1 and not Inertia2 then begin Prior:=true; exit; end;
   if Inertia2 and not Inertia1 then begin Prior:=false; exit; end;
 end;
end;

procedure IDA_Star;
var
 Sucess:boolean;
 CurrentState:StateType;
 H:integer;
 f:Text;

 procedure IDA_Push(State:StateType);
 begin
  if IDA.Top=MaxStack then
   Exit;
  inc(IDA.Top);
  IDA.Stack[IDA.Top]:=State;
 end;

 procedure IDA_Pop(var State:StateType);
 begin
  State:=IDA.Stack[IDA.Top];
  dec(IDA.Top);
 end;

 function IDA_Empty:boolean;
 begin
  IDA_Empty:=(IDA.Top=0);
 end;
\end{verbatim}

Trên đây là các thao tác ngăn xếp.

\begin{verbatim}
 procedure IDA_AddToHashTable(State:StateType);
 var
  h:integer;
  p:PHashTableEntry;
 begin
  h:=HashFunction(State);
  if HashTable^.Count[h]<MaxSubEntry then
  begin
    new(p);
    p^.State:=State;
    p^.Next:=HashTable^.FirstEntry[h];
    HashTable^.FirstEntry[h]:=p;
    inc(HashTable^.Count[h]);
  end
  else begin
    p:=HashTable^.FirstEntry[h];
    while p^.Next^.Next<>nil do
     p:=p^.Next;
    p^.Next^.State:=State;
    p^.Next^.Next:=HashTable^.FirstEntry[h];
    HashTable^.FirstEntry[h]:=p^.Next;
    p^.Next:=nil;
  end;
 end;

 function IDA_InHashTable(State:StateType):boolean;
 var
  h:integer;
  p:PHashTableEntry;
 begin
  h:=HashFunction(State);
  p:=HashTable^.FirstEntry[h];
  IDA_InHashTable:=true;
  while p<>nil do
  begin
   if SameState(p^.State,State) then
   begin
     if p^.State.g>State.g then
     begin
       p^.State.g:=State.g;
       IDA_InHashTable:=false;
\end{verbatim}

Nếu mục tìm thấy trong bảng có độ sâu lớn hơn, điều đó không có nghĩa là lần
này với độ sâu nhỏ hơn cũng vô nghiệm. Đúng ra nên cập nhật động cận dưới; ở
đây ta xử lý như chưa tìm thấy. Các chương sau sẽ cải tiến chỗ này.

\begin{verbatim}
     end;
     exit;
   end;
   p:=p^.Next;
  end;
  IDA_InHashTable:=false;
 end;
\end{verbatim}

Đó là các thao tác bảng băm.

\begin{verbatim}
 procedure IDA_AddNode(State:StateType);
 begin
  IDA_Push(State);
  inc(IDA.NodeCount);
  if IDA.NodeCount mod DispNode=0 then
    Writeln('NodeCount=',IDA.NodeCount);
  inc(IDA.TopLevelNodeCount);
  IDA_AddToHashTable(State);
 end;

 procedure IDA_Expand(State:StateType);
 var
  MoveCount:integer;
  MoveList:array[1..Maxx*Maxy*4] of MoveType;
  t:MoveType;
  i,j,Direction:integer;
  NewBoxPos, NewManPos:integer;
  NewState:StateType;
 begin
  MoveCount:=0;
  for i:=1 to MaxPosition do
   if GetBit(State.Boxes,i)>0 then
     for Direction:=0 to 3 do
     begin
      NewBoxPos:=i+DeltaPos[Direction];
      NewManPos:=i+DeltaPos[Opposite[Direction]];
      if GetBit(State.Boxes,NewBoxPos)>0 then continue;
      if GetBit(SokoMaze.Walls,NewBoxPos)>0 then continue;
      if GetBit(State.Boxes,NewManPos)>0 then continue;
      if GetBit(SokoMaze.Walls,NewManPos)>0 then continue;
      if CanReach(State,NewManPos) then
      begin
        DoMove(State,i,Direction,NewState);
        if CalcHeuristicFunction(NewState)=Infinite then continue;
        if CalcHeuristicFunction(NewState)+State.g>=IDA.PathLimit then continue;
        {Cot loi cua IDA*: gioi han do sau}
        if IDA_InHashTable(NewState) then continue;
        inc(MoveCount);
        MoveList[MoveCount].Position:=i;
        MoveList[MoveCount].Direction:=Direction;
      end;
     end;

  for i:=1 to MoveCount do
   for j:=i+1 to MoveCount do
    if Prior(State,MoveList[i],MoveList[j]) then
    {Dieu chinh thu tu cach day}
    begin
      t:=MoveList[j];
      MoveList[j]:=MoveList[i];
      MoveList[i]:=t;
    end;

  for i:=1 to MoveCount do
  {Lan luot xet moi phuong an di chuyen}
  begin
   DoMove(State,MoveList[i].Position,MoveList[i].Direction,NewState);
   inc(NewState.MoveCount);
   NewState.Move[NewState.MoveCount].Position:=MoveList[i].Position;
   NewState.Move[NewState.MoveCount].Direction:=MoveList[i].Direction;
   NewState.g:=State.g+1;
   IDA_AddNode(NewState);
  end;
 end;

 procedure IDA_Answer(State:StateType);
 var
  i:integer;
  x,y:integer;
 begin
  Writeln(f,'Solution Found in ', State.MoveCount,' Pushes');
  for i:=1 to State.Movecount do
  begin
    x:=State.Move[i].Position div SokoMaze.Y+1;
    y:=State.Move[i].Position mod SokoMaze.Y+1;
    Writeln(f, x,' ',y,' ',DirectionWords[State.Move[i].Direction]);
  end;
 end;

begin
 Writeln(VerStr);
 Writeln(Author);

 IDA.PathLimit:=CalcHeuristicFunction(IDA.StartState)-1;
 Sucess:=false;
 repeat
  inc(IDA.PathLimit);
  Writeln('Pathlimit=',IDA.PathLimit);
  IDA.TopLevelNodeCount:=0;
  IDA.Top:=0;
  IDA.StartState.g:=0;
  IDA_Push(IDA.StartState);
  repeat
   IDA_Pop(CurrentState);
   H:=CalcHeuristicFunction(CurrentState);
   if H=Infinite then continue;
   if Solution(CurrentState) then
     Sucess:=true
   else if IDA.PathLimit>=CurrentState.g+H then
     IDA_Expand(CurrentState);
  until Sucess or IDA_Empty or (IDA.NodeCount>MaxNode);
  Writeln('PathLimit ',IDA.PathLimit,' Finished. NodeCount=',IDA.NodeCount);
 until Sucess or (IDA.PathLimit>=MaxDepth) or (IDA.NodeCount>MaxNode);

 Assign(f,outfile);
 ReWrite(f);
 Writeln(f,VerStr);
 Writeln(f,Author);
 Writeln(f);

 if not Sucess then
   Writeln(f,'Cannot find a solution.')
 else
   IDA_Answer(CurrentState);

 Writeln('Node Count:',IDA.NodeCount);
 Writeln;
 close(f);
end;

begin
 Init;
 IDA_Star;
end.
\end{verbatim}

\section*{Ghi chú về bản trích xuất}
\addcontentsline{toc}{section}{Ghi chú về bản trích xuất}

Tệp PDF có tên ``incomplete'' và phần văn bản trích được bằng
\texttt{pdftotext -layout -enc UTF-8} kết thúc sau mã nguồn Baby-4 ở trên.
Các mục tiếp theo trong mục lục và các phụ lục chỉ xuất hiện dưới dạng tiêu
đề trong phần đầu của bản trích xuất.

\end{document}
